\section{\ENG{IMS} \ENG{Global} \ENG{Learning} \ENG{Consortium} (1997-σήμερα)}
\label{sec:ims_global}

Παράλληλα με το \ENG{ADL} και το \ENG{SCORM}, το \textit{\ENG{IMS} \ENG{Global} \ENG{Learning} \ENG{Consortium}} (αργότερα μετονομασμένο σε \textit{\ENG{1EdTech} \ENG{Consortium}}) αποτέλεσε έναν από τους πιο σημαντικούς οργανισμούς προτυποποίησης στον τομέα της εκπαιδευτικής τεχνολογίας. Το \ENG{IMS} εστίασε στην ανάπτυξη ανοιχτών προτύπων για τη διαλειτουργικότητα εκπαιδευτικών συστημάτων, με έμφαση στην ενσωμάτωση εργαλείων, την ανταλλαγή περιεχομένου και τη διαχείριση μαθησιακών πληροφοριών~\cite{ims2019lti,wiki_lti}.

\subsection{Ίδρυση και Όραμα του \ENG{IMS}}
\label{subsec:ims_foundation}

Το \ENG{IMS} \ENG{Global} \ENG{Learning} \ENG{Consortium} ιδρύθηκε το 1997 ως μια συνεργασία μεταξύ ακαδημαϊκών ιδρυμάτων, εταιρειών τεχνολογίας και εκπαιδευτικών οργανισμών. Η έδρα του βρισκόταν στο \ENG{EDUCAUSE}, μια μη κερδοσκοπική ένωση που υποστηρίζει την αξιοποίηση τεχνολογιών πληροφορικής στην ανώτατη εκπαίδευση~\cite{oneedtech2024lti}.

Ο όρος «\ENG{IMS}» αρχικά σήμαινε \textit{\ENG{Instructional} \ENG{Management} \ENG{Systems}}, αλλά με την πάροδο του χρόνου ο οργανισμός επέκτεινε το πεδίο του πέρα από τη διαχείριση και ο όρος έγινε απλώς ένα \ENG{brand} \ENG{name.} Το 2022, το \ENG{IMS} \ENG{Global} επισήμως μετονομάστηκε σε \ENG{1EdTech} \ENG{Consortium} για να αντικατοπτρίσει την ευρύτερη αποστολή του στην εκπαιδευτική τεχνολογία.

Το \ENG{IMS} λειτουργεί με βάση ένα μοντέλο μελών (\textit{\ENG{membership} \ENG{model}}), όπου οργανισμοί πληρώνουν συνδρομή για να συμμετάσχουν στην ανάπτυξη προτύπων, να έχουν πρόσβαση σε τεχνικά έγγραφα και να λάβουν πιστοποίηση συμμόρφωσης. Τα μέλη περιλαμβάνουν πανεπιστήμια, εταιρείες λογισμικού, κυβερνητικούς οργανισμούς και μη κερδοσκοπικές οργανώσεις.

\subsection{\ENG{Learning} \ENG{Tools} \ENG{Interoperability} (\ENG{LTI})}
\label{subsec:lti}

Το πιο σημαντικό πρότυπο που ανέπτυξε το \ENG{IMS} είναι το \textit{\ENG{Learning} \ENG{Tools} \ENG{Interoperability}} (\ENG{LTI}), το οποίο επιτρέπει την ενσωμάτωση εξωτερικών εκπαιδευτικών εργαλείων σε \ENG{LMS} χωρίς να απαιτείται \ENG{custom} \ENG{integration} για κάθε εργαλείο~\cite{ims2019overview,cetis2012lti}.

\subsubsection{Η Πρόκληση της Ενσωμάτωσης Εργαλείων}

Πριν από το \ENG{LTI}, η ενσωμάτωση εξωτερικών εργαλείων (π.χ. εργαλεία αξιολόγησης, συνεργατικές πλατφόρμες, εφαρμογές προσομοίωσης) σε ένα \ENG{LMS} απαιτούσε προσαρμοσμένο κώδικα για κάθε συνδυασμό εργαλείου-\ENG{LMS.} Αυτό οδηγούσε σε υψηλό κόστος ανάπτυξης και συντήρησης των \ENG{integrations}, αργή υιοθέτηση νέων εκπαιδευτικών εργαλείων, πολλαπλές συνδέσεις χρηστών (\ENG{multiple} \ENG{logins}) και αδυναμία μεταφοράς βαθμολογιών και δεδομένων προόδου μεταξύ συστημάτων.

Το \ENG{LTI} λύνει αυτό το πρόβλημα ορίζοντας ένα κοινό πρωτόκολλο επικοινωνίας μεταξύ του \ENG{LMS} (\textit{\ENG{Tool} \ENG{Consumer}} ή \textit{\ENG{Platform}}) και των εξωτερικών εργαλείων (\textit{\ENG{Tool} \ENG{Provider}} ή \textit{\ENG{Tool}}).

\subsubsection{\ENG{LTI} 1.0 και 1.1: Τα Θεμέλια (2010-2012)}

Το \ENG{LTI} 1.0 κυκλοφόρησε το 2010, ακολουθούμενο από το \ENG{LTI} 1.1 το 2012. Αυτές οι εκδόσεις παρείχαν βασική λειτουργικότητα. Το \ENG{Single} \ENG{Sign-On} (\ENG{SSO}) επιτρέπει στον εκπαιδευόμενο να συνδέεται μόνο στο \ENG{LMS}, ενώ το \ENG{LMS} δημιουργεί ένα ασφαλές αίτημα (\ENG{signed} με \ENG{OAuth} 1.0) που αυθεντικοποιεί τον χρήστη στο εργαλείο χωρίς νέα σύνδεση. Το \ENG{Context} \ENG{Information} μεταδίδει πληροφορίες (π.χ. το όνομα του μαθήματος, τον ρόλο του χρήστη—εκπαιδευόμενος, εκπαιδευτής, διαχειριστής) στο εργαλείο. Τα \ENG{Basic} \ENG{Outcomes} επιτρέπουν στο εργαλείο να στείλει πίσω στο \ENG{LMS} μια βαθμολογία (\ENG{score}) για την προσπάθεια του εκπαιδευόμενου.

Το \ENG{LTI} 1.1 χρησιμοποιεί ένα απλό \ENG{POST} \ENG{request} που περιέχει \ENG{parameters} όπως:

{\selectlanguage{english}\fontencoding{T1}\selectfont

\begin{verbatim}
user_id, roles, context_id, context_title, 
resource_link_id, lis_outcome_service_url, 
oauth_consumer_key, oauth_signature
\end{verbatim}
}

Το εργαλείο επαληθεύει την υπογραφή \ENG{OAuth} 1.0 για να διασφαλίσει ότι το αίτημα είναι αυθεντικό και στη συνέχεια παρέχει πρόσβαση στον χρήστη~\cite{moodle2024lti,blackboard2024lti}.

\subsubsection{\ENG{LTI} 1.3 και \ENG{LTI} \ENG{Advantage} (2017-2019)}

Το 2017, το \ENG{IMS} κυκλοφόρησε το \ENG{LTI} 1.3, μια σημαντική αναβάθμιση που αντικατέστησε το \ENG{OAuth} 1.0 με σύγχρονες μεθόδους ασφαλείας. Η αναβάθμιση στηρίχθηκε στη χρήση \ENG{OpenID} \ENG{Connect} (\ENG{OIDC}) για αυθεντικοποίηση και υποστήριξη \ENG{JSON} \ENG{Web} \ENG{Tokens} (\ENG{JWT}), στη χρήση \ENG{OAuth} 2.0 για \ENG{authorization} και σε \ENG{public/private} \ENG{key} \ENG{pairs} αντί για \ENG{shared} \ENG{secrets}, μειώνοντας τον κίνδυνο διαρροής διαπιστευτηρίων.

Το \ENG{LTI} \ENG{Advantage} είναι μια σειρά επεκτάσεων στο \ENG{LTI} 1.3 που προσθέτουν νέες δυνατότητες. Το \ENG{Deep} \ENG{Linking} επιτρέπει στους εκπαιδευτές να επιλέγουν συγκεκριμένο περιεχόμενο από το εργαλείο (π.χ. ένα συγκεκριμένο \ENG{quiz}) αντί να συνδέουν ολόκληρη την εφαρμογή. Οι \ENG{Assignment} \ENG{and} \ENG{Grade} \ENG{Services} (\ENG{AGS}) ενισχύουν τη μεταφορά βαθμολογιών και επιτρέπουν στο εργαλείο να δημιουργεί νέα στήλη βαθμολογίας (\ENG{grade} \ENG{column}) στο \ENG{LMS.} Οι \ENG{Names} \ENG{and} \ENG{Role} \ENG{Provisioning} \ENG{Services} (\ENG{NRPS}) δίνουν στο εργαλείο πρόσβαση στη λίστα χρηστών και ρόλων του μαθήματος, διευκολύνοντας ομαδικές δραστηριότητες και συνεργασία~\cite{d2l2024lti,unicon2024lti}.

\subsubsection{Αρχιτεκτονική και Ροή του \ENG{LTI}}

Η βασική ροή του \ENG{LTI} 1.3 ξεκινά όταν ο εκπαιδευόμενος κάνει κλικ σε ένα σύνδεσμο εργαλείου μέσα στο \ENG{LMS.} Το \ENG{LMS} εκκινεί μια \ENG{OpenID} \ENG{Connect} \ENG{authentication} \ENG{flow}, στέλνοντας ένα \textit{\ENG{login} \ENG{initiation} \ENG{request}} στο εργαλείο, και το εργαλείο ανακατευθύνει τον χρήστη πίσω στο \ENG{LMS} με ένα \textit{\ENG{authentication} \ENG{request}}. Στη συνέχεια το \ENG{LMS} δημιουργεί ένα \ENG{JWT} (\ENG{ID} \ENG{Token}) με πληροφορίες χρήστη και πλαισίου, το υπογράφει με το \ENG{private} \ENG{key} του και το στέλνει στο εργαλείο. Το εργαλείο επαληθεύει το \ENG{JWT} με το \ENG{public} \ENG{key} του \ENG{LMS} (από το \ENG{jwks}\_\ENG{uri} \ENG{endpoint}) και, αν η επαλήθευση είναι επιτυχής, παρέχει πρόσβαση στον χρήστη.

Αυτή η ροή εξασφαλίζει ότι μόνο εξουσιοδοτημένα \ENG{LMS} μπορούν να εκκινήσουν συνεδρίες και ότι τα δεδομένα χρήστη προστατεύονται κρυπτογραφικά.

\subsection{\ENG{Common} \ENG{Cartridge}}
\label{subsec:common_cartridge}

Ένα άλλο σημαντικό πρότυπο του \ENG{IMS} είναι το \textit{\ENG{Common} \ENG{Cartridge}} (\ENG{CC}), το οποίο ορίζει έναν τρόπο συσκευασίας και ανταλλαγής ψηφιακού εκπαιδευτικού περιεχομένου μεταξύ διαφορετικών \ENG{LMS}~\cite{ims2015cc,ims2015ccfaq}.

Το \ENG{Common} \ENG{Cartridge} είναι παρόμοιο με το \ENG{SCORM}, αλλά με σημαντικές διαφορές. Το \ENG{CC} μπορεί να περιλάβει όχι μόνο εκπαιδευτικά αντικείμενα (όπως το \ENG{SCORM}) αλλά και ολόκληρα μαθήματα με συζητήσεις, \ENG{forums}, \ENG{quizzes} και \ENG{assignments}, υποστηρίζει την ενσωμάτωση \ENG{LTI} \ENG{links} για αναφορά σε εξωτερικά εργαλεία, επιτρέπει τη συμπερίληψη \ENG{assessments} ορισμένων με το \ENG{IMS} \ENG{Question} \ENG{and} \ENG{Test} \ENG{Interoperability} (\ENG{QTI}) πρότυπο και χρησιμοποιεί \ENG{IMS} \ENG{metadata} \ENG{profile} (βασισμένο στο \ENG{IEEE} \ENG{LOM} και \ENG{Dublin} \ENG{Core})~\cite{coursearc2024cc}.

Η δομή ενός \ENG{Common} \ENG{Cartridge} περιλαμβάνει:

{\selectlanguage{english}\fontencoding{T1}\selectfont

\begin{verbatim}
imsmanifest.xml    <- Main manifest (declaration) file
resources/         <- Content folder (HTML, images, videos)
assessments/       <- QTI assessments
discussions/       <- Discussion topics
weblinks/          <- External links
\end{verbatim}
}

Το \ENG{Common} \ENG{Cartridge} υποστηρίζεται από μεγάλα \ENG{LMS} όπως το \ENG{Moodle}, το \ENG{Blackboard}, το \ENG{Canvas} και το \ENG{D2L} \ENG{Brightspace}~\cite{blackboard2024cc,moodle2024cc}.

\subsection{Άλλα Πρότυπα του \ENG{IMS/1EdTech}}

Πέρα από το \ENG{LTI} και το \ENG{Common} \ENG{Cartridge}, το \ENG{IMS/1EdTech} έχει αναπτύξει πολλά άλλα πρότυπα. Μεταξύ των βασικών προτύπων περιλαμβάνονται το \ENG{Question} \ENG{and} \ENG{Test} \ENG{Interoperability} (\ENG{QTI}) για τη δημιουργία και ανταλλαγή ερωτήσεων και τεστ, το \ENG{Learning} \ENG{Information} \ENG{Services} (\ENG{LIS}) για την ανταλλαγή πληροφοριών χρηστών, μαθημάτων και βαθμολογιών, το \ENG{Competencies} \ENG{and} \ENG{Academic} \ENG{Standards} \ENG{Exchange} (\ENG{CASE}) για μαθησιακούς στόχους και ικανότητες, τα \ENG{Open} \ENG{Badges} για ψηφιακά σήματα αναγνώρισης δεξιοτήτων και το \ENG{Caliper} \ENG{Analytics} για τη συλλογή και ανάλυση δεδομένων μαθησιακής αναλυτικής.

\subsection{Η Μετάβαση στο \ENG{1EdTech} (2022)}
\label{subsec:1edtech_rebrand}

Το 2022, το \ENG{IMS} \ENG{Global} ανακοίνωσε την επίσημη μετονομασία του σε \ENG{1EdTech} \ENG{Consortium}. Η αλλαγή αντικατοπτρίζει την επέκταση πέρα από την «\ENG{instructional} \ENG{management}» στην ευρύτερη εκπαιδευτική τεχνολογία, την έμφαση στην καινοτομία και την προσαρμογή στις σύγχρονες ανάγκες, καθώς και τη διεθνή εστίαση και συνεργασία με εκπαιδευτικούς οργανισμούς παγκοσμίως.

Η επωνυμία «\ENG{1EdTech}» υποδηλώνει «\ENG{One} \ENG{Education} \ENG{Technology}»—μια ενοποιημένη προσέγγιση στην προτυποποίηση. Όλα τα υπάρχοντα πρότυπα (\ENG{LTI}, \ENG{Common} \ENG{Cartridge}, \ENG{QTI} κ.λπ.) παραμένουν ενεργά και συνεχίζουν να εξελίσσονται υπό τη νέα επωνυμία.

\subsection{Επιρροή του \ENG{IMS/1EdTech} στην Εκπαίδευση}

Το \ENG{IMS} \ENG{Global} και το \ENG{1EdTech} έχουν διαδραματίσει κρίσιμο ρόλο στη διαμόρφωση του σύγχρονου οικοσυστήματος \ENG{E-Learning}. Η επίδραση αυτή φαίνεται στη διαλειτουργικότητα εργαλείων μέσω του \ENG{LTI}, που επιτρέπει την ενσωμάτωση χιλιάδων εκπαιδευτικών εργαλείων σε \ENG{LMS}, στη μείωση του κόστους ενσωμάτωσης και στην επιτάχυνση της υιοθέτησης νέων τεχνολογιών, στην ευελιξία επιλογής εργαλείων χωρίς δέσμευση σε έναν προμηθευτή και στην ενίσχυση της ασφάλειας μέσω σύγχρονων προτύπων όπως \ENG{LTI} 1.3, \ENG{OAuth} 2.0 και \ENG{OpenID} \ENG{Connect}.

Το \ENG{IMS/1EdTech} συνεχίζει να είναι ενεργό στην ανάπτυξη νέων προτύπων και στην προσαρμογή των υπαρχόντων για να ανταποκρίνονται στις εξελισσόμενες ανάγκες της εκπαίδευσης, συμπεριλαμβανομένης της τεχνητής νοημοσύνης, της εξατομικευμένης μάθησης και της διά βίου εκπαίδευσης.


% \ENG{Figure}: \ENG{IMS} \ENG{LTI} \ENG{Launch} \ENG{Flow}
\begin{figure}[H]
  \centering
  \begin{figure}[H]
\centering
\begin{chapterfigurebox}
\begin{tikzpicture}[
  every node/.style={font=\small},
  actor/.style={rectangle, draw, thick, minimum width=3.5cm, minimum height=1cm, align=center, fill=white},
  message/.style={-latex, thick},
  return/.style={-latex, thick, dashed}
]

% Actors
\node[actor, fill=imslightblue!30] at (0,9) (platform) {\textbf{\ENG{LMS/Platform}}\\(\ENG{Moodle}, \ENG{Canvas})};
\node[actor, fill=gray!20] at (6,9) (browser) {\textbf{\ENG{User Browser}}};
\node[actor, fill=imslightblue!50] at (12,9) (tool) {\textbf{\ENG{External Tool}}\\(\ENG{H5P}, \ENG{Virtual Lab})};

% Lifelines
\draw[dashed] (platform.south) -- (0,0);
\draw[dashed] (browser.south) -- (6,0);
\draw[dashed] (tool.south) -- (12,0);

% Step 1
\draw[message] (0,7.8) -- (6,7.8);
\node[above] at (3,7.8) {\ENG{1. User selects tool}};

% Step 2
\draw[message, imslightblue] (6,6.8) -- (12,6.8);
\node[above, align=center] at (9,6.8)
{\ENG{2. OIDC Login Initiation}\\
\texttt{\ENG{/oidc/login}}};

% Step 3
\node at (12,5.9) {\ENG{3. Tool prepares auth request}};

% Step 4
\draw[return, imslightblue] (12,5) -- (6,5);
\draw[message, imslightblue] (6,5) -- (0,5);
\node[above, align=center] at (3,5)
{\ENG{4. OAuth Authorization}\\
\texttt{\ENG{/oauth2/authorize}}};

% Step 5
\node at (0,4) {\ENG{5. Platform validates user}};

% Step 6
\draw[message, scormgreen] (0,3.1) -- (6,3.1);
\draw[message, scormgreen] (6,3.1) -- (12,3.1);
\node[above, align=center] at (9,3.1)
{\ENG{6. Launch request with ID Token (JWT)}};

% Step 7
\node at (12,2.2) {\ENG{7. Tool validates JWT}};

% Step 8
\draw[return, solutiongreen] (12,1.3) -- (6,1.3);
\node[above] at (9,1.3) {\ENG{8. Tool UI displayed}};

% Optional step
\draw[dashed, xapiorange, -latex] (12,0.5) -- (0,0.5);
\node[above] at (6,0.5)
{\ENG{9. Optional: Grade passback (AGS)}};

\end{tikzpicture}
\end{chapterfigurebox}

\vspace{0.4cm}

\begin{minipage}{0.45\textwidth}
\footnotesize
\textbf{\ENG{LTI} 1.3 \ENG{Security}}
\begin{itemize}
\item \ENG{OAuth} 2.0 + \ENG{OIDC}
\item \ENG{JWT} \ENG{signatures} (\ENG{RS256})
\item \ENG{Nonce} \ENG{validation}
\item \ENG{HTTPS} \ENG{communication}
\end{itemize}
\end{minipage}
\hfill
\begin{minipage}{0.45\textwidth}
\footnotesize
\textbf{\ENG{LTI} \ENG{Advantages}}
\begin{itemize}
\item \ENG{Single Sign-On}
\item \ENG{Standardized integration}
\item \ENG{Grade passback} (\ENG{AGS})
\item \ENG{Names} \& \ENG{Roles} (\ENG{NRPS})
\end{itemize}
\end{minipage}

\caption{\ENG{IMS LTI 1.3 Launch Flow using OAuth 2.0 and OpenID Connect}}
\label{fig:lti_launch_flow}

\end{figure}

  \label{fig:lti_flow}
\end{figure}

